\documentclass[11pt]{article}

\usepackage[margin=1in]{geometry}
\usepackage{amsmath,amssymb}
\usepackage{booktabs}
\usepackage{enumitem}
\usepackage{hyperref}
\usepackage{xcolor}

\hypersetup{
  colorlinks=true,
  linkcolor=blue!50!black,
  urlcolor=blue!50!black,
  citecolor=blue!50!black
}

\newcommand{\grip}{\texttt{grip}}
\newcommand{\Rfunc}[1]{\texttt{#1}}
\newcommand{\misf}{\mathrm{MISF}}
\newcommand{\RR}{\mathbb{R}}
\newcommand{\norm}[1]{\left\lVert #1 \right\rVert}
\newcommand{\eps}{\varepsilon}

\title{A MISF Edge-KK Layout with GRIP-Style Repulsion}
\author{Design note for the \grip{} package}
\date{May 19, 2026}

\begin{document}

\maketitle

\begin{abstract}
This document proposes a new layout algorithm, tentatively named
\Rfunc{grip.layout.misf.edge.kk()}, that merges two workflows currently used in
the \grip{} package: \Rfunc{grip.layout.weighted()} followed by
\Rfunc{grip.optimize.edge.kk.layout()}, or metric MDS followed by
\Rfunc{grip.optimize.edge.kk.layout()}. The revised specification is a staged
hybrid: weighted-MISF construction, heuristic insertion, explicit Armijo
objective refinement on each active MISF level, and a final pure or near-pure
\Rfunc{grip.optimize.edge.kk.layout()} polish. During MISF refinement, the
objective separates true original-edge stress from temporary graph-geodesic
metric-neighbor stress, log repulsion, and quadratic anchors. The goal is to
guide the layout toward edge-isometric shape throughout the filtration without
calling temporary metric-neighbor constraints ``edge-KK'' and without using
local-temperature updates in the version-0 objective refinement.
\end{abstract}

\tableofcontents

\section{Motivation}

The current high-quality stretched-isometric workflow is usually one of:
\begin{align*}
  Z_0 &\leftarrow \Rfunc{grip.layout.weighted()}(G),\\
  Z   &\leftarrow \Rfunc{grip.optimize.edge.kk.layout()}(G, Z_0),
\end{align*}
or
\begin{align*}
  Z_0 &\leftarrow \Rfunc{grip.metric.mds.layout()}(G),\\
  Z   &\leftarrow \Rfunc{grip.optimize.edge.kk.layout()}(G, Z_0).
\end{align*}
This split is useful because each piece contributes something different.
\Rfunc{grip.layout.weighted()} provides a multiscale GRIP construction with
repulsion and good global stabilization. \Rfunc{grip.optimize.edge.kk.layout()}
then improves edge-length fidelity under the edge-restricted KK objective.

The separation also has a limitation: the edge-KK objective only acts after the
major multiscale placement decisions have already been made. If the weighted
GRIP layout has chosen a folded or compressed basin, edge-KK can repair edge
lengths locally but may not easily move the embedding into a better global
configuration. A more direct algorithm should use the edge-KK objective during
the MISF construction itself while retaining the repulsive stabilization that
makes GRIP useful on sparse graphs.

\section{Relationship to Existing Functions}

The proposed function should not replace any existing function in place.
It should be introduced beside the current code paths as an experimental
algorithm:
\[
  \Rfunc{grip.layout.misf.edge.kk()}.
\]

\begin{table}[h]
\centering
\begin{tabular}{lll}
\toprule
Function & Current role & Proposed relationship \\
\midrule
\Rfunc{grip.layout.weighted()} &
Weighted GRIP layout with FR-like final refinement &
Baseline and source of MISF/repulsion mechanics \\
\Rfunc{grip.layout.weighted.nd()} &
Dimension-general weighted-GRIP backend under development &
Future infrastructure model for arbitrary dimension support \\
\Rfunc{grip.optimize.edge.kk.layout()} &
Edge-KK local repair optimizer &
Source of the edge-stress objective and scale policies \\
\Rfunc{grip.metric.mds.layout()} &
Metric-MDS initializer for edge-KK workflows &
Baseline initializer and comparison method \\
\Rfunc{grip.layout.misf.edge.kk()} &
Not yet implemented &
New integrated MISF edge-KK plus repulsion solver \\
\bottomrule
\end{tabular}
\caption{Relationship between the proposed method and existing \grip{} tools.}
\end{table}

The design principle is conservative: keep existing public behavior intact,
reuse tested components where possible, and build the new method as an
opt-in algorithm with its own tests and diagnostics.

\section{Notation}

Let
\[
  G=(V,E), \qquad |V|=n,
\]
be an undirected graph. Each edge \(e=(i,j)\in E\) has a positive target length
\[
  \ell_{ij} > 0.
\]
In current \grip{} terminology, these target lengths come from edge weights or
normalized edge weights. The embedding is
\[
  Z = (z_1,\ldots,z_n)^\top \in \RR^{n\times d},
\]
where \(d \ge 2\).

For an edge \((i,j)\), define the regularized Euclidean length
\[
  r_{ij}(Z) = \sqrt{\norm{z_i-z_j}^2 + \eps_{\mathrm{edge}}^2}.
\]
Let \(k_{ij} > 0\) be the edge-KK stiffness. In the existing edge-KK optimizer,
stiffnesses may be uniform, density-weighted, or distance-power weighted.

The weighted-MISF hierarchy is represented by an ordered vertex array
\[
  \pi = (\pi_1,\ldots,\pi_n),
\]
called \texttt{misf.order} in this document, and level sizes
\[
  m_L < m_{L-1} < \cdots < m_0=n.
\]
The active set at level \(q\) is defined exactly as the prefix
\[
  A_q = \{\pi_1,\ldots,\pi_{m_q}\}.
\]
Thus the MISF hierarchy is a nested sequence
\[
  A_L \subset A_{L-1} \subset \cdots \subset A_0 = V,
\]
where \(A_L\) is the coarsest level and \(A_0\) is the full graph. When moving
from level \(q+1\) to level \(q\), the newly inserted vertices are
\[
  N_q = A_q \setminus A_{q+1}.
\]
The active-level optimizer moves only vertices in \(A_q\). Vertices outside
\(A_q\) do not appear in edge, metric, repulsion, or anchor constraints at level
\(q\). Previously active vertices \(A_{q+1}\) and newly inserted vertices
\(N_q\) may both move during Armijo refinement.

\section{Target Objective}

The version-0 method is a staged hybrid. Weighted-MISF construction and
heuristic insertion define the active set and starting coordinates for each
level. The level itself is refined by an explicit Armijo objective, not by a
GRIP local-temperature force update. For a level \(q\), the objective is
\[
  \mathcal{E}_q(Z_{A_q})
  =
  \mathcal{E}_{\mathrm{edge},q}(Z_{A_q})
  +
  \rho_q \mathcal{E}_{\mathrm{metric},q}(Z_{A_q})
  +
  \lambda_q \mathcal{E}_{\mathrm{rep},q}(Z_{A_q})
  +
  \mathcal{E}_{\mathrm{anchor},q}(Z_{A_q}),
\]
where all terms involve only vertices in \(A_q\). The temporary metric-neighbor
constraints are auxiliary graph-geodesic constraints. They are not called
edge-KK constraints. Only the original-edge term is edge-KK.

\subsection{True Original-Edge Stress}

The true edge term is restricted to original graph edges whose endpoints are
active at level \(q\):
\[
  E_{\mathrm{active},q}
  =
  \{(i,j)\in E : i,j\in A_q\}.
\]
Its stress is
\[
  \mathcal{E}_{\mathrm{edge},q}(Z)
  =
  \frac{1}{2}
  \sum_{(i,j)\in E_{\mathrm{active},q}}
  k_{ij}^{(q)}
  \left(
    r_{ij}(Z) - s_q \ell_{ij}
  \right)^2.
\]
Here \(k_{ij}^{(q)}>0\) is the edge stiffness, \(s_q>0\) is the active-level
edge scale, and \(\ell_{ij}\) is the normalized original edge target length.

The gradient contribution for a true active edge \((i,j)\) is
\[
  \nabla_{z_i}
  \mathcal{E}_{ij}
  =
  k_{ij}^{(q)}
  \frac{r_{ij}(Z)-s_q\ell_{ij}}{r_{ij}(Z)}
  (z_i-z_j),
\]
with the opposite sign for \(z_j\). This is the edge-stress gradient used by
\Rfunc{grip.optimize.edge.kk.layout()}.

\subsection{Auxiliary Metric-Neighbor Stress}

Coarse MISF levels may have too few original edges among active vertices.
Version 0 therefore allows a separate auxiliary metric-neighbor constraint set
\[
  M_q
  =
  \{(i,j,\delta_{ij}) : i,j\in A_q,\ i<j,\ \delta_{ij}=d_G(i,j)\}.
\]
Here \(d_G(i,j)\) is the weighted graph-geodesic distance. The set \(M_q\) is
constructed from weighted metric-neighbor caches, deduplicated as unordered
pairs, and kept disjoint from \(E_{\mathrm{active},q}\). If an original active
edge appears in the metric-neighbor cache, it remains only in
\(E_{\mathrm{active},q}\) and is removed from \(M_q\).

The auxiliary metric stress is
\[
  \mathcal{E}_{\mathrm{metric},q}(Z)
  =
  \frac{1}{2}
  \sum_{(i,j,\delta_{ij})\in M_q}
  h_{ij}^{(q)}
  \left(
    r_{ij}(Z) - s_q^{\mathrm{metric}}\delta_{ij}
  \right)^2.
\]
The coefficient \(\rho_q\) outside this term controls how much these temporary
metric constraints affect the level. The schedule must satisfy
\[
  \rho_0 = 0
\]
at the full graph level by default, so the final full-level objective is not
held to auxiliary metric-neighbor constraints.

\subsection{Log Repulsion}

The repulsive term preserves the stabilizing role currently played by GRIP:
\[
  \mathcal{E}_{\mathrm{rep},q}(Z)
  =
  -
  \sum_{(i,j)\in P_q}
  \omega_{ij}^{(q)}
  \log\left(
    \sqrt{\norm{z_i-z_j}^2+\eps_{\mathrm{rep}}^2}
  \right).
\]
The pair set \(P_q\) contains only active vertices. It is exact for small active
sets and deterministic sampled for larger active sets. Crucially, \(P_q\) and
its weights are fixed before an Armijo iteration begins and remain fixed across
all trial steps in that iteration. The objective being tested by Armijo is
therefore stable during line search.

The negative gradient contribution on \(z_i\) is
\[
  -\nabla_{z_i}
  \lambda_q \mathcal{E}_{\mathrm{rep},q}
  =
  \lambda_q \omega_{ij}^{(q)}
  \frac{z_i-z_j}{\norm{z_i-z_j}^2+\eps_{\mathrm{rep}}^2}.
\]
Pair weights must be normalized so the repulsion magnitude does not grow like
\(|A_q|^2\). A default normalization is
\[
  \sum_{(i,j)\in P_q}\omega_{ij}^{(q)} = |A_q|,
\]
or equivalently a fixed average repulsion mass per active vertex.

\subsection{Quadratic Anchors}

After insertion and before level refinement, store the active coordinates as
\[
  \tilde{Z}^{(q)} = Z_{A_q}^{\mathrm{post\ insertion}}.
\]
The anchor term is
\[
  \mathcal{E}_{\mathrm{anchor},q}(Z)
  =
  \frac{1}{2}
  \sum_{i\in A_q}
  a_i^{(q)}
  \norm{z_i-\tilde{z}_i^{(q)}}^2.
\]
Anchor weights should be strongest immediately after insertion, decay during
later levels, and be off for the final pure edge-KK polish. Anchors are a
level-transition stabilizer, not part of the final edge-KK objective.

\section{Scale Policy}

Scale handling is intentionally restricted during MISF refinement. While
repulsion or auxiliary metric stress is active, the edge scale should be either
identity or fixed. Profiled scale should not be used during repulsive MISF
refinement because repulsion changes the state being profiled and can cause
level-dependent scale drift.

Version 0 supports the following MISF-level scale policies:
\begin{description}[leftmargin=2em]
  \item[\texttt{identity}] Use \(s_q=1\) for all MISF levels.
  \item[\texttt{global\_fixed}] Fit one scale before MISF refinement begins, or
  use a user-supplied scale, and keep it fixed for all MISF levels.
  \item[\texttt{level\_fixed}] Fit \(s_q\) once immediately before refining
  level \(q\), then keep it fixed throughout that level's Armijo iterations.
  This is diagnostic and should not be the first default.
  \item[\texttt{user}] Use a user-supplied fixed scale.
\end{description}

Profiled scale remains available for the final
\Rfunc{grip.optimize.edge.kk.layout()} polish and for diagnostics with
\(\lambda_q=\rho_q=0\). For profiled edge-only polish, the least-squares
estimate is
\[
  s
  =
  \frac{
    \sum_{(i,j)\in E}
    k_{ij} r_{ij}(Z) \ell_{ij}
  }{
    \sum_{(i,j)\in E}
    k_{ij} \ell_{ij}^2
  },
\]
with the finite-value safeguards already used by
\Rfunc{grip.optimize.edge.kk.layout()}.

\section{Level-Specific Constraint Sets}

For each level \(q\), construct four disjoint or separately tracked objects:
\[
  E_{\mathrm{active},q}, \qquad M_q, \qquad P_q, \qquad B_q.
\]
Here \(B_q\) denotes the anchor target and anchor weights, not graph edges.

\subsection{Original Active Edges}

\[
  E_{\mathrm{active},q}
  =
  \{(i,j)\in E : i,j\in A_q\}.
\]
These are the only true edge-KK constraints during MISF refinement.

\subsection{Temporary Metric-Neighbor Constraints}

\[
  M_q
  =
  \{(i,j,\delta_{ij}) : i,j\in A_q,\ i<j,\ \delta_{ij}=d_G(i,j)\}.
\]
The set is built from weighted metric-neighbor caches, deduplicated, and
explicitly separated from original active edges. These constraints are
temporary graph-geodesic stabilizers. They are controlled by \(\rho_q\) and
decay to zero at the full graph level.

\subsection{Active Repulsion Pairs}

\[
  P_q \subseteq \{(i,j): i,j\in A_q,\ i<j\}.
\]
Use exact active pairs below a threshold such as
\(|A_q|\le n_{\mathrm{rep,exact}}\). Above the threshold, use a deterministic
sample keyed by level, seed, and active vertex order. The sampled set and its
normalized weights are fixed for all Armijo trial evaluations at the iteration
where they are used.

\subsection{Anchor Targets}

Anchor targets are the coordinates immediately after insertion and before
level refinement:
\[
  B_q =
  \left\{
    (i,\tilde{z}_i^{(q)},a_i^{(q)}) : i\in A_q
  \right\}.
\]
Both previously active vertices and newly inserted vertices can have anchors.
Weights may be larger for \(N_q\) than for \(A_{q+1}\) at early iterations, but
this is a tunable default.

\section{Algorithm Overview}

The algorithm has six conceptual phases:
\begin{enumerate}[leftmargin=2em]
  \item Validate graph, weights, dimension, and tuning parameters.
  \item Build a weighted MISF hierarchy.
  \item Initialize the coarsest active set and heuristically insert vertices at
  each subsequent level.
  \item At each active level, build separate true-edge, auxiliary metric,
  repulsion, and anchor objects.
  \item Refine all active vertices by Armijo descent on the explicit
  level objective.
  \item Run a final pure or near-pure
  \Rfunc{grip.optimize.edge.kk.layout()} polish.
\end{enumerate}

\subsection{Inputs}

A proposed R signature is:
\begin{verbatim}
grip.layout.misf.edge.kk <- function(
  edges = NULL,
  n = NULL,
  adj_list = NULL,
  weight_list = NULL,
  edge_weights = NULL,
  dim = 3L,
  num_init = 24L,
  num_nbrs = 20L,
  rounds = 80L,
  final_rounds = 160L,
  auxiliary_metric = TRUE,
  repulsion_mode = c("hybrid", "active_exact", "active_sampled",
                     "none"),
  repulsion_factor = 1.0,
  repulsion_schedule = c("coarse_to_fine", "constant", "off_full"),
  metric_factor = 0.25,
  metric_schedule = c("coarse_to_fine", "constant", "off_full"),
  anchor_factor = 0.1,
  anchor_schedule = c("coarse_to_fine", "constant", "off_full"),
  stiffness_method = c("density", "uniform", "distance_power"),
  stiffness_transform = c("identity", "sqrt", "log"),
  density_mix_schedule = c(0, 0.5, 1),
  misf_scale_mode = c("identity", "global_fixed", "level_fixed", "user"),
  final_scale_mode = c("profiled", "identity", "fixed_initial", "user"),
  scale = NULL,
  edge_length_epsilon = 1e-8,
  repulsion_epsilon = 1e-8,
  initial_step = 1.0,
  step_shrink = 0.5,
  armijo_factor = 1e-4,
  grad_tol = 1e-8,
  min_step = 1e-8,
  final_edge_kk_polish = TRUE,
  final_edge_kk_iter = 50L,
  seed = 6L,
  return_trace = FALSE,
  disconnected = c("components", "error")
)
\end{verbatim}

This signature is intentionally verbose for design discussion. The first
implementation should expose fewer controls and keep experimental controls
internal until their behavior is understood.

\section{Detailed Algorithm}

\subsection{Step 1: Preprocessing}

Validate and normalize the graph exactly as weighted-GRIP does:
\begin{enumerate}[leftmargin=2em]
  \item Accept either edge matrix input or adjacency-list input.
  \item Validate positive edge weights.
  \item Normalize edge lengths using the existing weighted layout conventions.
  \item Build adjacency and weight lists.
  \item For disconnected graphs, either error or solve components separately
  and pack them, matching existing weighted-GRIP behavior.
\end{enumerate}

Let the normalized target edge length be \(\ell_{ij}\).

\subsection{Step 2: Build Weighted MISF}

Reuse the weighted MISF construction from the weighted-GRIP/ND infrastructure.
The output should include:
\begin{itemize}[leftmargin=2em]
  \item the ordered vertex array \(\pi=\texttt{misf.order}\),
  \item MISF level sizes \(m_q=\texttt{misf.level\_size[q]}\),
  \item weighted metric-neighbor caches,
  \item insertion anchors,
  \item active-level metadata for tracing.
\end{itemize}

The MISF should be built from weighted graph distances, not unweighted graph
distance, because the edge-KK target lengths are weighted. At every later level
the active set is exactly the prefix \(A_q=\{\pi_1,\ldots,\pi_{m_q}\}\), and
newly inserted vertices are \(N_q=A_q\setminus A_{q+1}\).

\subsection{Step 3: Coarsest-Level Initialization}

Let \(A_L\) be the coarsest active set.
Possible initialization modes:
\begin{description}[leftmargin=2em]
  \item[\texttt{random\_centered}] current GRIP-style random initial placement,
  centered at the origin.
  \item[\texttt{metric\_mds}] metric-MDS on the coarsest weighted metric cache.
  \item[\texttt{simplex}] deterministic near-simplex initialization when
  \(|A_L|\le d+1\).
\end{description}

The recommended initial default is \texttt{metric\_mds} when cheap and
\texttt{random\_centered} otherwise. For parity with existing weighted-GRIP
diagnostics, \texttt{random\_centered} should remain available.

\subsection{Step 4: Insert New Vertices}

When moving from \(A_{q+1}\) to \(A_q\), each new vertex \(v\in A_q\setminus
A_{q+1}\) is assigned an initial coordinate before refinement.

The insertion objective should be a small local edge/metric KK problem:
\[
  \min_{z_v\in\RR^d}
  \frac{1}{2}
  \sum_{a\in B_q(v)}
  b_{va}
  \left(
    \sqrt{\norm{z_v-z_a}^2+\eps_{\mathrm{edge}}^2}
    -
    d_G(v,a)
  \right)^2
  +
  \lambda_{\mathrm{ins}}
  \mathcal{R}_{\mathrm{local}}(z_v),
\]
where \(B_q(v)\) is an anchor set among already-active vertices and
\(d_G(v,a)\) is the weighted graph distance. The local repulsion term is
optional and should usually be weak.

Practical insertion modes:
\begin{description}[leftmargin=2em]
  \item[\texttt{barycenter}] weighted barycenter of anchors.
  \item[\texttt{least\_squares}] distance-fitting least-squares insertion.
  \item[\texttt{local\_edge\_kk}] a few Armijo steps on the insertion objective.
\end{description}

The recommended default is:
\[
  \texttt{least\_squares} + \texttt{local\_edge\_kk\_polish}.
\]

\subsection{Step 5: Active-Level Armijo Refinement}

For each level \(q\), run a fixed-budget Armijo optimizer on the explicit
objective \(\mathcal{E}_q\). The optimized variables are only
\(\{z_i:i\in A_q\}\). Vertices outside \(A_q\) are absent from the objective and
do not move. Both previously active vertices in \(A_{q+1}\) and newly inserted
vertices in \(N_q\) may move.

Before the first Armijo iteration at level \(q\):
\begin{enumerate}[leftmargin=2em]
  \item construct \(E_{\mathrm{active},q}\), \(M_q\), \(P_q\), and \(B_q\);
  \item compute \(\rho_q\), \(\lambda_q\), anchor weights, stiffnesses, and the
  fixed MISF-level scale;
  \item normalize repulsion pair weights so the total repulsion mass is
  \(O(|A_q|)\), not \(O(|A_q|^2)\);
  \item store anchor targets from the post-insertion coordinates.
\end{enumerate}

Each Armijo iteration evaluates
\[
  G_q = \nabla_{Z_{A_q}}\mathcal{E}_q(Z_{A_q})
\]
and proposes
\[
  Z_{A_q}' = Z_{A_q} - \eta G_q.
\]
The proposal is accepted if
\[
  \mathcal{E}_q(Z_{A_q}')
  \le
  \mathcal{E}_q(Z_{A_q})
  -
  c\eta\norm{G_q}_F^2.
\]
Otherwise the step is shrunk, \(\eta\leftarrow\gamma\eta\), until the Armijo
condition is met or \(\eta<\eta_{\min}\).

After an accepted step, recenter the active coordinates \(Z_{A_q}\). Recenter
only active coordinates for the current component; inactive vertices do not
participate.

\subsubsection{Stopping Criteria}

Stop level \(q\) when any of the following holds:
\begin{enumerate}[leftmargin=2em]
  \item the fixed iteration budget \(T_q\) is reached;
  \item \(\norm{G_q}_F \le \tau_{\mathrm{grad}}\);
  \item no Armijo step above \(\eta_{\min}\) is accepted;
  \item relative objective decrease over a small patience window falls below a
  tolerance.
\end{enumerate}

\subsubsection{Recommended First Implementation}

The recommended first implementation should use Armijo descent for every MISF
refinement level. Local-temperature updates are deliberately excluded from
version 0 objective refinement because they obscure the objective being
optimized. They may be revisited later as an acceleration heuristic only after
the Armijo objective path is validated.

\section{Level Schedules}

Version 0 has three non-edge schedules: auxiliary metric strength \(\rho_q\),
repulsion strength \(\lambda_q\), and anchor weights \(a_i^{(q)}\). All three
must decay toward the final graph level so the final polish can be interpreted
as pure or near-pure edge-KK.

\subsection{Auxiliary Metric Schedule}

The metric-neighbor coefficient \(\rho_q\) controls temporary graph-geodesic
constraints:
\[
  \rho_L \ge \rho_{L-1} \ge \cdots \ge \rho_0 = 0.
\]
The default should use positive \(\rho_q\) on coarse levels and switch it off
at the full graph level. A simple default is
\[
  \rho_q = \rho_{\max}
  \left(
    \frac{q}{L}
  \right)^{\gamma_\rho},
\]
with \(\rho_0=0\) enforced exactly.

\subsection{Repulsion Schedule}

The repulsion coefficient \(\lambda_q\) should vary by level. A useful default
is:
\[
  \lambda_q
  =
  \lambda_0
  \left(
    \frac{|A_q|}{|V|}
  \right)^\gamma
  +
  \lambda_{\min},
\]
or a simple piecewise schedule:
\[
  \lambda_q =
  \begin{cases}
    \lambda_{\mathrm{coarse}}, & \text{coarse levels},\\
    \lambda_{\mathrm{mid}},    & \text{intermediate levels},\\
    \lambda_{\mathrm{final}},  & \text{full graph}.
  \end{cases}
\]

The default should probably use stronger repulsion during coarse and insertion
levels and weaker repulsion during final full-level refinement:
\[
  \lambda_{\mathrm{coarse}}
  >
  \lambda_{\mathrm{mid}}
  >
  \lambda_{\mathrm{final}}
  \ge 0.
\]
By default, \(\lambda_{\mathrm{final}}\) should be zero or near zero before the
final \Rfunc{grip.optimize.edge.kk.layout()} polish. If nonzero full-level
repulsion is used, the final polish should be described as near-pure rather
than pure edge-KK.

This lets repulsion protect the topology early without preventing final
edge-length accuracy.

\subsection{Anchor Schedule}

Anchor weights are stored as \(a_i^{(q)}\). They should stabilize the transition
from insertion to refinement, then decay:
\[
  a_i^{(q,t)}
  =
  a_{i,0}^{(q)}
  \left(1-\frac{t}{T_q}\right)^{\gamma_a},
\]
where \(t\) is the Armijo iteration index within level \(q\). Anchors are off
for final polish. A tunable default is to use a larger initial anchor for newly
inserted vertices \(N_q\) than for already active vertices \(A_{q+1}\).

\subsection{Repulsion Pair Normalization}

For exact repulsion, \(|P_q|=|A_q|(|A_q|-1)/2\), which would make the raw log
repulsion scale quadratically with active count. Version 0 therefore normalizes
weights so
\[
  \sum_{(i,j)\in P_q} \omega_{ij}^{(q)} = |A_q|.
\]
For sampled repulsion, use the same total mass after sampling. The pair set and
weights are fixed during Armijo trial evaluations.

\section{Stiffness Schedule}

The existing edge-KK optimizer supports a density-mix continuation schedule.
The MISF version should generalize this to levels:
\[
  k_{ij}^{(q,t)}
  =
  (1-\mu_t)k_{ij}^{\mathrm{density}}
  +
  \mu_t k_{ij}^{\mathrm{uniform}},
\]
where \(\mu_t\in[0,1]\) is a continuation value. The current default schedule
\[
  (0,0.25,0.5,0.75,1)
\]
is probably too expensive if applied at every MISF level. A smaller default
such as
\[
  (0,0.5,1)
\]
or a level-dependent schedule should be considered.

One practical design is:
\begin{itemize}[leftmargin=2em]
  \item coarse levels: density-weighted stiffnesses to stabilize irregular
  edge-length distributions;
  \item intermediate levels: mixed stiffnesses;
  \item final level: uniform or nearly uniform stiffnesses for faithful
  edge-isometry.
\end{itemize}

This stiffness schedule applies separately to true edge constraints and
auxiliary metric constraints. Metric-neighbor stiffnesses should be treated as
temporary stabilizers and scaled by \(\rho_q\); they should not be reported as
edge-KK stiffnesses.

\section{Trace Schema}

The trace should make the staged hybrid explicit. Each accepted or attempted
Armijo iteration should record at least:
\begin{itemize}[leftmargin=2em]
  \item level index \(q\);
  \item active count \(|A_q|\);
  \item inserted count \(|N_q|\);
  \item counts of true edges \(|E_{\mathrm{active},q}|\), metric constraints
  \(|M_q|\), and repulsion pairs \(|P_q|\);
  \item \(\rho_q\), \(\lambda_q\), anchor summary weights, and repulsion weight
  normalization;
  \item scale mode and scale value;
  \item total objective and separate edge, metric, repulsion, and anchor
  energies;
  \item gradient norm;
  \item accepted step size and number of line-search shrinkages;
  \item true-edge relative RMSE;
  \item auxiliary metric relative RMSE when \(|M_q|>0\);
  \item elapsed time for the iteration or level.
\end{itemize}

\section{Pseudocode}

\subsection{High-Level Solver}

\begin{verbatim}
function grip.layout.misf.edge.kk(graph, weights, dim, controls):
    validated = validate_weighted_graph(graph, weights, dim)
    components = split_components_if_requested(validated)
    layouts = []

    for component in components:
        hierarchy = build_weighted_misf(component, controls)
        coords = initialize_coarsest_level(hierarchy.top, dim, controls)
        previous_active = empty_set

        for level from top_level down to full_graph_level:
            active = prefix(hierarchy.misf_order,
                            hierarchy.level_size[level])
            newly_inserted = active \ previous_active

            if level is not top_level:
                coords = insert_new_vertices(coords, hierarchy, level, controls)

            anchors = capture_post_insertion_anchor_targets(coords, active)
            true_edges = original_edges_with_both_endpoints_in(active)
            metric_constraints = metric_neighbor_constraints(
                hierarchy, active, excluding = true_edges
            )
            repulsion_pairs = build_repulsion_pairs(hierarchy, level, controls)
            schedules = resolve_level_schedules(level, controls)

            coords = armijo_refine_active_level(
                coords,
                active,
                newly_inserted,
                true_edges,
                metric_constraints,
                repulsion_pairs,
                anchors,
                schedules,
                controls
            )

            record_trace_if_requested()
            previous_active = active

        if controls.final_edge_kk_polish:
            coords = grip.optimize.edge.kk.layout(coords = coords, ...)

        layouts.append(coords)

    return pack_components(layouts)
\end{verbatim}

\subsection{Level Optimizer}

\begin{verbatim}
function armijo_refine_active_level(coords, active, newly_inserted,
                                    true_edges, metric_constraints,
                                    repulsion_pairs, anchors, schedules,
                                    controls):
    for stage in stiffness_schedule:
        edge_stiffness = build_edge_stiffness(true_edges, stage)
        metric_stiffness = build_metric_stiffness(metric_constraints, stage)
        scale = fixed_misf_scale(coords, true_edges, controls)

        for iter in 1:rounds_for_level:
            state = evaluate_level_objective(
                coords,
                active,
                true_edges,
                edge_stiffness,
                metric_constraints,
                metric_stiffness,
                repulsion_pairs,
                anchors,
                schedules,
                scale
            )

            if state.gradient_norm < grad_tol:
                break

            proposal = armijo_line_search(coords, active, state, controls)
            if proposal.accepted:
                coords[active, ] = proposal.coords[active, ]
                recenter_active_or_component(coords, active)
            else:
                break

            record_level_trace(state, proposal)

    return coords
\end{verbatim}

\section{Implementation Plan}

\subsection{Phase 0: Design and Baselines}

\begin{enumerate}[leftmargin=2em]
  \item Keep this design document as the reference.
  \item Create benchmark scripts comparing:
  \begin{itemize}
    \item \Rfunc{grip.layout.weighted()} followed by
    \Rfunc{grip.optimize.edge.kk.layout()};
    \item \Rfunc{grip.metric.mds.layout()} followed by
    \Rfunc{grip.optimize.edge.kk.layout()};
    \item the future \Rfunc{grip.layout.misf.edge.kk()}.
  \end{itemize}
  \item Use existing fixtures: paths, cycles, grids, nonuniform grids, lattices,
  quadforms, trees, and irregular manifold-like graphs.
\end{enumerate}

\subsection{Phase 1: R Prototype}

Implement a clear R prototype before writing compiled code:
\begin{enumerate}[leftmargin=2em]
  \item Build or reuse a weighted MISF object.
  \item At each level, define \(A_q\) as the prefix of \texttt{misf.order} and
  \(N_q=A_q\setminus A_{q+1}\).
  \item Heuristically insert \(N_q\), then capture anchor targets.
  \item Build \(E_{\mathrm{active},q}\), \(M_q\), \(P_q\), and anchor objects as
  separate data structures.
  \item Use the existing R edge-KK gradient function for true edge terms.
  \item Add R implementations of auxiliary metric stress, log repulsion, and
  quadratic anchors.
  \item Use Armijo descent for every active-level refinement.
  \item Return coordinates, trace, and diagnostics.
\end{enumerate}

The R prototype will be slower, but it will make the objective and trace schema
easy to verify.

\subsection{Phase 2: Minimal Compiled Backend}

Once the R prototype is correct, implement the hot loop in C++:
\begin{enumerate}[leftmargin=2em]
  \item Add a compiled state evaluator for true edge stress, auxiliary metric
  stress, log repulsion, and quadratic anchors.
  \item Support arbitrary dimension.
  \item Support exact active-pair repulsion below a threshold.
  \item Support deterministic sampled repulsion above the threshold.
  \item Support identity, global fixed, level fixed, and user scale modes during
  MISF refinement.
  \item Keep profiled scale out of repulsive MISF refinement and reserve it for
  final edge-KK polish or edge-only diagnostics.
  \item Return per-stage trace rows with counts, schedules, term energies,
  gradient norms, accepted step sizes, RMSEs, and elapsed time.
\end{enumerate}

This phase could reuse ideas from:
\begin{itemize}[leftmargin=2em]
  \item \texttt{src/edge\_isometric\_gkk\_rcpp.cpp},
  \item \texttt{src/edge\_repulsive\_unfolding\_rcpp.cpp},
  \item \texttt{src/DrawGraphND.cpp}.
\end{itemize}

\subsection{Phase 3: MISF Integration}

Integrate the compiled evaluator into a MISF engine:
\begin{enumerate}[leftmargin=2em]
  \item Reuse weighted MISF construction.
  \item Reuse weighted insertion anchors for heuristic insertion, not as hidden
  objective terms.
  \item Add level-wise construction of \(E_{\mathrm{active},q}\), \(M_q\),
  \(P_q\), and \(B_q\).
  \item Ensure sampled \(P_q\) is deterministic and fixed during Armijo trials.
  \item Add trace frames compatible with the staged hybrid trace schema.
\end{enumerate}

At this phase the method becomes a true MISF edge-KK algorithm rather than a
post-hoc optimizer.

\subsection{Phase 4: Final Polish and Public Wrapper}

Add the exported R wrapper:
\[
  \Rfunc{grip.layout.misf.edge.kk()}.
\]
The wrapper should:
\begin{itemize}[leftmargin=2em]
  \item use familiar weighted-GRIP input conventions;
  \item support \(\dim\ge 2\);
  \item optionally return trace and diagnostics;
  \item run a final pure or near-pure edge-KK polish by default;
  \item clearly mark the function as experimental until benchmarks stabilize.
\end{itemize}

\subsection{Phase 5: Testing and Reports}

Add tests in layers:
\begin{enumerate}[leftmargin=2em]
  \item validation tests for inputs and dimensions;
  \item finite-output tests in dimensions 2, 3, and 4;
  \item edge-stress decrease tests on small graphs;
  \item full objective decrease tests for each active-level Armijo refinement;
  \item equivalence tests where repulsion is off and one level is used, checking
  agreement with \Rfunc{grip.optimize.edge.kk.layout()};
  \item tests that \(M_q\) is deduplicated and disjoint from
  \(E_{\mathrm{active},q}\);
  \item tests that sampled \(P_q\) is deterministic and fixed during line
  search;
  \item comparison reports against the two existing pipelines;
  \item gflow-backed scientific validation for grids, nonuniform grids, and
  quadform examples.
\end{enumerate}

\section{Tuning Strategy for Auxiliary Metric and Repulsion Schedules}

The two most important regularization schedules are:
\[
  \rho_q \quad \text{for auxiliary metric-neighbor stress},
  \qquad
  \lambda_q \quad \text{for log repulsion}.
\]
These schedules should be tuned deliberately because they control the central
trade-off of the method. Large \(\rho_q\) and \(\lambda_q\) help the MISF
construction choose stable global basins, but they can also pull the final
embedding away from edge-isometric fidelity. Small values make the method closer
to pure edge-KK, but may allow coarse-level ambiguity, collapse, or fold-over.

\subsection{Recommended Tuning Philosophy}

Version 0 should not use a fully adaptive controller. The first release should
use a small global schedule family with graph-size and level normalization, plus
trace diagnostics that make failures interpretable. Adaptive modifiers should
be added later only when a diagnostic has a clear meaning.

The recommended progression is:
\begin{enumerate}[leftmargin=2em]
  \item tune a robust global default schedule family on a broad fixture panel;
  \item keep pair-mass normalization and level-fraction normalization active
  from the beginning;
  \item add simple adaptive guards only after the global default is stable;
  \item keep all adaptive decisions visible in the trace.
\end{enumerate}

This keeps the method auditable. A global default answers the question
``what algorithm are we running?'' Adaptive guards should answer narrower
questions such as ``was this active edge graph too sparse for pure edge
constraints at this level?''

\subsection{Schedule Parameterization}

Use a low-dimensional schedule family rather than arbitrary per-level values.
Let
\[
  \theta_q =
  \begin{cases}
    q/L, & L>0,\\
    0, & L=0,
  \end{cases}
\]
so that \(\theta_q=1\) at the coarsest level and \(\theta_q=0\) at the full
graph level. A default auxiliary metric schedule can be
\[
  \rho_q =
  \rho_{\max}\theta_q^{\gamma_\rho},
  \qquad
  \rho_0=0.
\]
A default repulsion schedule can be
\[
  \lambda_q =
  \lambda_{\min}
  +
  \lambda_{\max}\theta_q^{\gamma_\lambda},
\]
with \(\lambda_{\min}=0\) for pure final edge-KK polish, or a small nonzero
\(\lambda_{\min}\) only when the final full-level refinement is intentionally
near-pure rather than pure edge-KK.

The first tuning grid should be small:
\[
  \rho_{\max}\in\{0,0.05,0.1,0.25,0.5,1.0\},
  \qquad
  \gamma_\rho\in\{0.5,1,2\},
\]
\[
  \lambda_{\max}\in\{0,0.05,0.1,0.25,0.5,1.0\},
  \qquad
  \gamma_\lambda\in\{0.5,1,2\}.
\]
The actual numerical scale may need adjustment after the objective terms are
implemented because edge stress, metric stress, and log repulsion can have
different natural magnitudes. The trace must therefore record term energies and
gradient norms, not only final layout quality.

\subsection{Fixture Panel}

The tuning panel should include graph families that stress different failure
modes:
\begin{itemize}[leftmargin=2em]
  \item paths and cycles, to detect over-regularization and scale drift;
  \item flat 2D grids, to detect fold-over and edge-length bias;
  \item nonuniform 2D grids, to test variable edge lengths;
  \item flat 3D lattices embedded in 3D and projected to 2D, to test
  cross-dimensional behavior;
  \item nonuniform 3D grids, to test irregular weighted geometry;
  \item quadform-derived 2D and 3D examples, to test local metric fidelity;
  \item trees and branching graphs, to test sparse active edge sets;
  \item irregular manifold-like graphs, such as annulus, torus, shell, sphere,
  and pair-of-pants examples where available;
  \item graphs with bottlenecks or density variation, to test whether
  repulsion and auxiliary metrics prevent collapse without tearing.
\end{itemize}

Each fixture should be run across several seeds unless the initialization is
fully deterministic. Seed stability is a first-class output, not a secondary
plot.

\subsection{Baselines}

Every tuning report should compare the same graph against:
\begin{enumerate}[leftmargin=2em]
  \item \Rfunc{grip.layout.weighted()} followed by
  \Rfunc{grip.optimize.edge.kk.layout()};
  \item \Rfunc{grip.metric.mds.layout()} followed by
  \Rfunc{grip.optimize.edge.kk.layout()};
  \item \Rfunc{grip.layout.misf.edge.kk()} with \(\rho_q=\lambda_q=0\) where
  this is meaningful;
  \item \Rfunc{grip.layout.misf.edge.kk()} over the candidate schedule grid.
\end{enumerate}

The zero-regularization MISF run is not expected to be best. It is a diagnostic
baseline that shows what the MISF machinery contributes without auxiliary
metric or repulsive stabilization.

\subsection{Evaluation Metrics}

The tuning objective should be multi-criteria. A candidate schedule should not
win solely because it improves one scalar stress while producing an unusable
shape. Record at least:
\begin{itemize}[leftmargin=2em]
  \item final true-edge relative RMSE;
  \item final true-edge signed bias;
  \item graph-geodesic distortion on sampled vertex pairs;
  \item auxiliary metric RMSE at levels where \(M_q\) is active;
  \item local angle or quadform preservation where ground-truth geometry exists;
  \item Procrustes RMSE to known coordinates for synthetic fixtures;
  \item fold-over or cell-orientation failures for grid-like fixtures;
  \item minimum pair distance and low-percentile pair distances;
  \item seed-to-seed variability of final metrics;
  \item number of Armijo iterations, rejected trial steps, and early stops;
  \item runtime and peak memory where measurable.
\end{itemize}

The trace-level metrics should include, per level:
\begin{itemize}[leftmargin=2em]
  \item edge, metric, repulsion, and anchor energy fractions;
  \item edge and metric gradient-norm fractions;
  \item \(|E_{\mathrm{active},q}|\), \(|M_q|\), and \(|P_q|\);
  \item active edge graph connected-component count;
  \item repulsion pair normalization mass;
  \item accepted step-size distribution.
\end{itemize}

\subsection{Selecting a Global Default}

The first global default should be selected as a Pareto compromise, not as the
minimizer of a single metric. A candidate default should satisfy:
\begin{enumerate}[leftmargin=2em]
  \item no systematic fold-over failures on grid-like fixtures;
  \item final edge RMSE competitive with the two current pipelines;
  \item better or comparable seed stability than the two current pipelines;
  \item no large graph family where auxiliary metric stress remains dominant at
  late levels;
  \item acceptable runtime relative to
  \Rfunc{grip.layout.weighted()} followed by
  \Rfunc{grip.optimize.edge.kk.layout()}.
\end{enumerate}

If no single schedule is robust, prefer a conservative default with weaker
regularization and expose stronger presets for difficult sparse or folding-prone
families.

\subsection{Adaptive Guards for Later Versions}

After a global default is established, adaptive guards can adjust \(\rho_q\) and
\(\lambda_q\) using interpretable diagnostics. Recommended guards are:
\begin{description}[leftmargin=2em]
  \item[Active-edge sparsity guard] If the graph induced by
  \(E_{\mathrm{active},q}\) has too many connected components or too small an
  average active degree, increase \(\rho_q\) for that level.
  \item[Collapse guard] If a low percentile of active pair distances falls below
  a scale-adjusted threshold, increase \(\lambda_q\) or delay its decay.
  \item[Repulsion-dominance guard] If repulsion gradient norm or energy remains
  too large relative to true edge stress late in the solve, decrease
  \(\lambda_q\).
  \item[Auxiliary-metric-dominance guard] If metric-neighbor stress dominates
  true edge stress at late levels, reduce \(\rho_q\) faster.
  \item[Armijo-failure guard] If line search repeatedly fails at a level, reduce
  the active regularization coefficients before reducing the edge objective.
\end{description}

Each adaptive change must be recorded in the trace with the diagnostic that
triggered it. Adaptive behavior should be optional until it is benchmarked
against the global default.

\subsection{Tuning Report Outputs}

The tuning report should produce:
\begin{itemize}[leftmargin=2em]
  \item a per-run metrics table with graph, seed, schedule parameters, final
  metrics, and runtime;
  \item a per-level trace summary table showing term energies and gradient
  fractions;
  \item Pareto plots of edge RMSE versus fold-over rate, edge RMSE versus
  geodesic distortion, and edge RMSE versus runtime;
  \item seed-stability plots for the best candidate schedules;
  \item a short recommended default and a list of fixtures where it fails.
\end{itemize}

This report should become the main evidence for whether \(\rho_q\) and
\(\lambda_q\) are fixed defaults, preset-specific defaults, or adaptive
quantities.

\section{Expected Advantages}

The integrated MISF edge-KK method should offer:
\begin{itemize}[leftmargin=2em]
  \item better edge-length fidelity than weighted-GRIP alone;
  \item better global basin selection than post-hoc edge-KK alone;
  \item less dependence on metric-MDS as an initializer;
  \item a natural path to high-dimensional edge-isometric embeddings;
  \item a direct framework for comparing repulsion schedules and edge-KK
  stiffness schedules.
\end{itemize}

\section{Risks and Open Questions}

\subsection{Repulsion Can Fight Edge Isometry}

Repulsion is essential for avoiding collapse and fold-over, but it can also
stretch edges away from their desired lengths. The solution should be a schedule
that uses repulsion strongly during coarse construction and weakly during final
edge-KK polish.

\subsection{Coarse Levels May Need Non-Edge Constraints}

Pure edge constraints among MISF vertices may be too sparse at coarse levels.
Hybrid metric constraints are likely necessary. The key question is how quickly
to remove or downweight them so the final result remains an edge-KK layout.

\subsection{Armijo Cost}

Armijo descent gives a clean objective and is required for version 0, but it
may be slower than the current GRIP local-temperature refinement. The first
implementation should prioritize objective clarity over speed. Local
temperature should remain out of the version-0 refinement path and may be
reconsidered only after the Armijo formulation is validated.

\subsection{Scale Drift Across Levels}

Profiled scale can change from level to level and can interact badly with
repulsion. For that reason, profiled scale is excluded from MISF refinement
while repulsion or auxiliary metric stress is active. The implementation should
record scale in the trace and support identity, global fixed, and level fixed
scale as diagnostic modes.

\subsection{Benchmark Definition}

The method should not be judged only by edge stress. It should be evaluated by:
\begin{itemize}[leftmargin=2em]
  \item edge relative RMSE,
  \item local angle or quadform preservation where available,
  \item graph-geodesic distance distortion,
  \item visual fold-over diagnostics,
  \item stability across seeds,
  \item runtime and memory use.
\end{itemize}

\section{Recommended First Milestone}

The first useful milestone should be deliberately narrow:
\begin{quote}
Implement an R prototype of a one-component MISF edge-KK layout in dimensions
\(d\ge2\), using hybrid coarse constraints, exact active repulsion for small
graphs, fixed or identity MISF-level scale, Armijo descent, and a final optional
call to
\Rfunc{grip.optimize.edge.kk.layout()}.
\end{quote}

Success criteria:
\begin{enumerate}[leftmargin=2em]
  \item finite deterministic layouts on paths, cycles, grids, and trees;
  \item objective decrease at each active level when using Armijo;
  \item when MISF has only one level and repulsion is disabled, results match
  \Rfunc{grip.optimize.edge.kk.layout()} within numerical tolerance;
  \item on canonical fixtures, the integrated method matches or improves the
  final edge relative RMSE of \Rfunc{grip.layout.weighted()} followed by
  \Rfunc{grip.optimize.edge.kk.layout()}.
\end{enumerate}

\section{Conclusion}

\Rfunc{grip.layout.misf.edge.kk()} is a natural extension of the current
\grip{} layout ecosystem. It should be viewed as a new staged hybrid algorithm
rather than a replacement for \Rfunc{grip.layout.weighted()}. The conceptual
merge is: keep the weighted-MISF hierarchy and insertion machinery, refine each
active level with an explicit Armijo objective containing true edge stress,
temporary auxiliary metric stress, log repulsion, and quadratic anchors, and
finish with a pure or near-pure edge-KK polish.

If implemented side by side, the method gives \grip{} a principled route toward
stretched isometric embeddings whose active-level construction is already
biased toward edge fidelity, while still using temporary metric and repulsive
regularization to choose a stable global basin.

\end{document}
